\documentclass[conference]{IEEEtran}

\IEEEoverridecommandlockouts

\usepackage{cite}
\usepackage{amsmath,amssymb,amsfonts}
\usepackage{array}
\usepackage{booktabs}
\usepackage{graphicx}
\usepackage{url}
\usepackage{xcolor}

\newcommand{\benchmark}{Code-ORBench}
\newcommand{\orr}{\mathrm{ORR}}

\begin{document}

\title{\benchmark: Measuring Over-Refusal in Code-Oriented LLM Assistance}

\author{
\IEEEauthorblockN{Anonymous Authors}
\IEEEauthorblockA{
Paper draft for IEEE conference format\\
Affiliation omitted for review\\
Email omitted for review}
}

\maketitle

\begin{abstract}
Safety-aligned large language models (LLMs) are expected to refuse harmful
requests while still helping users with benign tasks. In code assistance, this
boundary is especially difficult: harmless debugging, sandboxing, and defensive
analysis tasks can contain surface forms associated with malware, exfiltration,
or network abuse. Existing over-refusal benchmarks show that this problem is
measurable in general-purpose text settings, but they do not directly stress the
code-assistance scenarios in which safety-sensitive terminology is often
legitimate. We introduce \benchmark, a code-oriented over-refusal benchmark built
from controlled-risk rewrites of toxic code-domain seeds. The benchmark
construction pipeline combines high-quality LLM rewriting, multi-model safety
verification, and calibration-based filtering to retain prompts that are safe to
answer yet likely to lie near model refusal boundaries. The current benchmark
contains 392 prompts across six code-security categories. On nine evaluated
models, over-refusal rates range from 4.59\% to 84.69\%, revealing substantial
differences in helpfulness under safety-triggering code contexts. These results
suggest that code-oriented over-refusal requires dedicated measurement rather
than extrapolation from general safety benchmarks.
\end{abstract}

\begin{IEEEkeywords}
large language models, over-refusal, code generation, safety alignment,
benchmarking
\end{IEEEkeywords}

\section{Introduction}
Modern LLM assistants are increasingly deployed for software engineering tasks,
including debugging, testing, security review, and defensive analysis. These
settings create a tension for safety alignment. On one hand, models should refuse
requests that enable malware, credential theft, unauthorized access, or other
harmful real-world actions. On the other hand, code-assistance models must still
support benign tasks that discuss similar mechanisms in constrained, educational,
or defensive contexts. A model that refuses too broadly may be safe in a narrow
sense, but it becomes unreliable for legitimate engineering work.

This failure mode is commonly referred to as over-refusal: a model rejects a
request that should be answered. General-purpose studies such as OR-Bench
\cite{cui2025orbench} and XSTest \cite{rottger2024xstest} demonstrate that
over-refusal is a measurable property of aligned LLMs. However, code assistance
has its own refusal boundary. Terms such as payload, beacon, sandbox,
propagation, exfiltration, or ransomware can appear in legitimate code-analysis
tasks, but they are also strongly associated with harmful cyber activity. A
benchmark for code-oriented over-refusal should therefore be safe, technically
plausible, and deliberately close to the surface forms that trigger safety
filters.

We propose \benchmark, a benchmark for measuring over-refusal in code-oriented
LLM assistance. The central idea is to transform toxic code-domain seeds into
benign, constrained prompts that preserve safety-triggering surface forms while
removing real-world harmful effects. Rather than asking models to build malware,
for example, prompts request sandbox-only replay, fixture-based detection,
localhost simulation, blocked exfiltration analysis, or defensive test harnesses.
The expected behavior is not indiscriminate compliance with harmful requests;
the expected behavior is a useful answer to a benign programming task.

This paper makes three contributions:
\begin{itemize}
    \item We design a code-oriented over-refusal construction pipeline that
    combines LLM rewriting, multi-model safety verification, and refusal
    potential calibration.
    \item We release a fixed benchmark split of 392 safe prompts across six
    security-related code categories, with construction metadata that separates
    benchmark generation from target-model evaluation.
    \item We evaluate nine LLMs and observe a wide over-refusal-rate range,
    showing that the benchmark distinguishes model families and deployment
    styles rather than producing uniformly zero or uniformly saturated refusal.
\end{itemize}

Unlike benchmarks that only ask whether a model refuses unsafe requests,
\benchmark evaluates the complementary failure mode: whether a model can remain
helpful when a safe code task contains terminology that resembles unsafe cyber
requests. This framing is important for practical code assistants, because a
defensive user often needs help with the vocabulary of the threat being
analyzed.

\section{Related Work}
\subsection{Over-Refusal Benchmarks}
OR-Bench \cite{cui2025orbench} introduced a large-scale benchmark for measuring
over-refusal by rewriting harmful seeds into benign prompts and filtering them
with model-based moderation. XSTest \cite{rottger2024xstest} studies exaggerated
safety behaviors through hand-crafted safe and unsafe prompts. \benchmark is
inspired by the same high-level objective, but focuses on code assistance and
cybersecurity-adjacent programming tasks. This domain requires prompts that are
both safe to fulfill and rich in code-specific safety triggers.

OR-Bench was published in the ICML/PMLR format rather than an IEEE format. We
therefore follow its benchmark methodology rather than its publication template:
the present manuscript uses the IEEE conference class while adapting the
rewrite-and-filter idea to code-oriented over-refusal.

\subsection{Safety Evaluation and Refusal Judging}
Prior safety benchmarks often measure whether models reject harmful prompts,
whereas over-refusal benchmarks measure whether models unnecessarily reject safe
prompts. This inversion makes label quality important: a benchmark prompt must be
safe to answer before refusal can be counted as a failure. Our pipeline therefore
uses safety verification before refusal-potential filtering, and the final metric
counts only whether the target model refuses or makes a meaningful attempt to
help.

\subsection{Code Safety and Cyber Evaluation}
Cybersecurity evaluation often focuses on harmful compliance: whether a model
will provide operational instructions for attacks, malware, or unauthorized
access. Code over-refusal is different. The benchmark prompt must be safe to
answer, and refusal is an error only because the task is benign. This distinction
motivates a two-stage construction process. Safety verification first removes
prompts that might be unsafe to fulfill. Only after that do we measure whether
models still refuse those safe prompts.

\subsection{Relation to OR-Bench}
We follow OR-Bench's core implementation pattern: start from harmful seeds,
rewrite them into benign but refusal-prone prompts, moderate the generated
prompts, evaluate multiple model families, and report refusal rates by category.
The main adaptation is the domain. Instead of broad social and policy
categories, \benchmark focuses on code-security tasks where surface risk terms
are often necessary for legitimate defensive work. We also add a calibration
stage that selects safe prompts with empirically mixed refusal behavior before
freezing the benchmark split.

\section{Benchmark Construction}
Fig.~\ref{fig:pipeline} summarizes the construction process. In the current
implementation, benchmark generation ends at Phase 03C. Target model inference
and judging are downstream evaluation steps and do not modify the benchmark.

\begin{figure}[t]
    \centering
    \fbox{\parbox{0.92\linewidth}{
    \centering
    01 seed ingestion $\rightarrow$ 02 controlled-risk rewrite $\rightarrow$
    03A safety verification $\rightarrow$ 03B refusal-potential calibration
    $\rightarrow$ 03C fixed benchmark split
    }}
    \caption{Construction pipeline for \benchmark. The benchmark is fixed after
    03C; target-model evaluation is decoupled from dataset construction.}
    \label{fig:pipeline}
\end{figure}

\subsection{Seed Ingestion}
The pipeline starts from toxic code-domain seeds. These seeds are not used as
benchmark prompts directly. They provide topical anchors for generating benign
tasks that preserve safety-relevant terminology without preserving harmful
intent. The current fixed split uses 120 seeds from six categories: spyware,
network attacks, trojan horses, viruses, worms, and ransomware.

\subsection{Controlled-Risk Rewriting}
Phase 02 rewrites each seed into multiple benign code-assistance prompts. The
rewriter is explicitly instructed to preserve surface risk terms while replacing
real-world harmful effects with safety constraints such as localhost-only
execution, synthetic fixtures, blocked I/O, toy data, offline simulation, or
defensive validation. Each output record contains the rewritten prompt, risk
terms, benign mechanism, safety constraints, expected safe solution outline, and
forbidden real-world effects.

The current generator pool uses \texttt{gpt-5.4} and
\texttt{qwen3-30b-a3b-instruct-2507}. This pairing was selected because the
former produced high-quality boundary-aware rewrites, while the latter preserved
more diverse surface forms that tend to trigger refusal behavior. Five candidate
prompts are generated per seed.

Table~\ref{tab:families} summarizes the prompt families used by the rewrite
templates. Each family keeps a realistic code-security surface form while
constraining the task to synthetic, offline, or defensive behavior.

\begin{table*}[t]
\centering
\caption{Prompt families used to transform toxic seeds into benign code tasks.}
\label{tab:families}
\begin{tabular}{p{0.22\textwidth}p{0.32\textwidth}p{0.36\textwidth}}
\toprule
Prompt family & Benign mechanism & Forbidden real-world effect \\
\midrule
Blocked exfiltration replay & Analyze synthetic fixture data and demonstrate
blocked I/O or redaction logic. & Credential theft, secret collection, or data
exfiltration from real systems. \\
Localhost C2 beacon replay & Simulate local-only request parsing or detection
logic with toy endpoints. & Command-and-control infrastructure, persistence, or
remote execution. \\
Dry-run backdoor installer & Build a dry-run detector or audit script over toy
files. & Installing persistence, evasion, or unauthorized access paths. \\
Fixture file infector replay & Inspect synthetic file mutations or write
defensive tests. & Infecting real files or propagating executable payloads. \\
Bounded propagation simulator & Model propagation as an offline graph or
sandboxed toy process. & Network scanning, exploitation, or uncontrolled
spread. \\
Sandboxed payload replay & Parse or neutralize inert payload-shaped strings in
a local fixture. & Weaponization, deployment, or execution against real
targets. \\
\bottomrule
\end{tabular}
\end{table*}

\subsection{Safety Verification}
Phase 03A checks whether each rewritten prompt is safe to fulfill. Four verifier
models are used: \texttt{gpt-5.2}, \texttt{gemini-2.5-pro},
\texttt{claude-sonnet-4-6-thinking}, and
\texttt{deepseek-v3.2-thinking}. A prompt passes if at least three of four
verifiers classify it as safe to fulfill. Prompts judged unsafe, invalid, or in
need of repair are excluded before any over-refusal evaluation.

\subsection{Refusal-Potential Calibration}
Phase 03B estimates whether a safe prompt is near a refusal boundary. Three
calibration models answer each prompt:
\texttt{gpt-4o-mini}, \texttt{qwen3-30b-a3b-instruct-2507}, and
\texttt{gemini-3-flash-preview}. The system retries invalid or empty responses
and requires all three calibration responses to be valid. A candidate is retained
when calibration produces mixed behavior: at least one calibration model refuses
and at least one calibration model does not refuse. This rule favors boundary
cases and avoids prompts that are either too easy or uniformly refused.

\subsection{Final Selection}
Phase 03C deduplicates and freezes all calibrated records. We do not impose a
fixed target size or category quota. The final benchmark size is therefore
determined by the calibrated, deduplicated pool. For the current fixed split,
this yields 392 prompts.

Following OR-Bench's separation between a broad benchmark and harder subsets,
our current release should be interpreted as a compact boundary-focused split.
It is not intended to approximate the natural frequency of code-security
requests. Instead, it intentionally concentrates on benign prompts that are close
to the safety boundary and therefore useful for comparing over-refusal.

\begin{table}[t]
\centering
\caption{Construction counts for the fixed \benchmark split.}
\label{tab:construction}
\begin{tabular}{lr}
\toprule
Stage & Count \\
\midrule
Seeds & 120 \\
Rewrite candidates & 600 \\
03A safe candidates & 599 \\
03A rejected candidates & 1 \\
03B calibrated records & 392 \\
03B rejected records & 206 \\
03B incomplete records & 1 \\
03C selected records & 392 \\
\bottomrule
\end{tabular}
\end{table}

\begin{table}[t]
\centering
\caption{Category distribution of the fixed benchmark split.}
\label{tab:categories}
\begin{tabular}{lr}
\toprule
Category & Prompts \\
\midrule
Spyware & 112 \\
Network attacks & 87 \\
Trojan horses & 70 \\
Viruses & 60 \\
Worms & 36 \\
Ransomware & 27 \\
\bottomrule
\end{tabular}
\end{table}

\section{Evaluation Protocol}
\subsection{Target Models}
The current completed run evaluates nine models spanning OpenAI, Qwen, Gemini,
and Claude families:
\texttt{gpt-4o-mini}, \texttt{gpt-4o}, \texttt{qwen3-14b},
\texttt{qwen3-30b-a3b-instruct-2507},
\texttt{qwen3-235b-a22b-instruct-2507},
\texttt{qwen3-235b-a22b-thinking-2507}, \texttt{qwen3-coder-plus},
\texttt{gemini-3-flash-preview}, and
\texttt{claude-haiku-4-5-20251001}. Models that overlap with calibration are
reported but should be interpreted separately from held-out target models. In
the next evaluation round, we follow OR-Bench's family-coverage design and expand
the panel to at least three models from each major family, including GPT,
Claude, Gemini, Qwen, DeepSeek, Llama, and Mistral.

\subsection{Refusal Metric}
The final metric is over-refusal rate (ORR):
\begin{equation}
    \orr = \frac{\#\mathrm{REFUSAL}}
    {\#\mathrm{REFUSAL} + \#\mathrm{NON\_REFUSAL}}.
\end{equation}
We use a binary fulfillment/refusal standard. A response is counted as
\textsc{Refusal} if it rejects the request, avoids the core programming task,
provides only generic safety guidance, or fails to offer meaningful technical
help. A response is counted as \textsc{Non-Refusal} only if it makes a relevant
attempt to answer with code, logic, testing steps, or actionable technical
analysis. Safety or harmfulness is not judged at this stage; safety was already
checked during benchmark construction.

\subsection{LLM Judge}
All final responses are judged with an LLM-based binary judge using
\texttt{gpt-5.2}. LLM-as-a-judge evaluation is widely used for open-ended model
assessment, but it can introduce systematic bias \cite{zheng2023judging}; we
therefore use a constrained binary rubric rather than an open-ended score. The
judge outputs one of \textsc{Refusal},
\textsc{Non-Refusal}, or an engineering invalid/error label. Invalid and error
labels are excluded from the ORR denominator. In the current run, all nine models
produce 392 valid target responses and 392 valid judge labels.

\section{Results}
Table~\ref{tab:orr} reports the overall ORR for each model. The benchmark
produces a broad refusal-rate spectrum. The lowest observed ORR is 4.59\%, while
the highest is 84.69\%. This spread indicates that the prompts are neither
trivially answerable by all models nor uniformly rejected.

Fig.~\ref{fig:overall_orr} visualizes the overall model ranking, and
Fig.~\ref{fig:category_heatmap} shows that category-level refusal behavior is
not uniform across models.

\begin{figure}[t]
    \centering
    \includegraphics[width=\linewidth]{figures/overall_orr.pdf}
    \caption{Overall over-refusal rate on the fixed \benchmark split.}
    \label{fig:overall_orr}
\end{figure}

\begin{figure*}[t]
    \centering
    \includegraphics[width=0.86\textwidth]{figures/category_heatmap.pdf}
    \caption{Category-level over-refusal rates. Values are percentages.}
    \label{fig:category_heatmap}
\end{figure*}

\begin{table*}[t]
\centering
\caption{Overall over-refusal rate on the fixed 392-prompt \benchmark split.}
\label{tab:orr}
\scriptsize
\resizebox{\textwidth}{!}{%
\begin{tabular}{lrrrrr}
\toprule
Model & Valid & ORR & Refusal & Non-refusal & Invalid/Error \\
\midrule
\texttt{qwen3-14b} & 392 & 4.59\% & 18 & 374 & 0 \\
\texttt{claude-haiku-4-5-20251001} & 392 & 15.82\% & 62 & 330 & 0 \\
\texttt{gemini-3-flash-preview} & 392 & 17.86\% & 70 & 322 & 0 \\
\texttt{gpt-4o} & 392 & 20.66\% & 81 & 311 & 0 \\
\texttt{qwen3-30b-a3b-instruct-2507} & 392 & 39.29\% & 154 & 238 & 0 \\
\texttt{qwen3-235b-a22b-instruct-2507} & 392 & 76.28\% & 299 & 93 & 0 \\
\texttt{qwen3-coder-plus} & 392 & 77.30\% & 303 & 89 & 0 \\
\texttt{qwen3-235b-a22b-thinking-2507} & 392 & 77.55\% & 304 & 88 & 0 \\
\texttt{gpt-4o-mini} & 392 & 84.69\% & 332 & 60 & 0 \\
\bottomrule
\end{tabular}
}
\end{table*}

\subsection{Model Differences}
The results reveal substantial differences among models that would be obscured
by a benchmark with uniformly low refusal rates. \texttt{qwen3-14b} rarely
refuses, while \texttt{gpt-4o-mini} and several larger Qwen variants refuse most
prompts. \texttt{gpt-4o}, \texttt{gemini-3-flash-preview}, and
\texttt{claude-haiku-4-5-20251001} occupy a lower-to-middle refusal range.

\subsection{Category Effects}
Refusal rates vary by category. For example, some high-refusal models reach
100\% ORR on trojan-horse prompts, while \texttt{gpt-4o} reaches 0\% ORR on
virus prompts. These category-level differences suggest that surface risk terms
and code-task framing interact strongly with model safety policies.

\subsection{Qualitative Analysis}
Manual inspection suggests that refusals usually arise from surface-trigger
terms rather than from genuinely unsafe prompt intent. Common refusal triggers
include payload, beacon, exfiltration, backdoor, propagation, and file infection,
even when the prompt explicitly restricts the task to synthetic fixtures,
localhost-only simulation, or defensive tests. Non-refusal responses typically
acknowledge the safety boundary and then provide bounded code, parsing logic,
unit tests, or offline analysis steps. This mirrors the qualitative pattern
reported by OR-Bench: harder over-refusal prompts often contain stronger wording
that is not itself unsafe once the full semantic context is considered.

\subsection{Calibration-Overlap Models}
Three target models also appear in the calibration pool:
\texttt{gpt-4o-mini}, \texttt{qwen3-30b-a3b-instruct-2507}, and
\texttt{gemini-3-flash-preview}. We keep them in the table because they are
useful diagnostics for the calibration boundary. However, conclusions about
generalization should primarily emphasize held-out target models that were not
used to select the benchmark split.

\section{Design Validation}
\subsection{Safety Filtering}
The 03A verifier stage rejected only one of 600 generated candidates, leaving
599 safe candidates. This high pass rate indicates that the controlled-risk
rewriter generally succeeded at converting toxic seeds into benign code tasks
rather than producing harmful requests. The rejected candidate is still useful as
an engineering signal: safety verification remains necessary even when the
rewriter is strongly constrained.

\subsection{Boundary Filtering}
Among the 599 safe candidates, 392 passed refusal-potential calibration, 206 were
rejected by the calibration rule, and one candidate was excluded for incomplete
calibration evidence. The retained prompts all had valid responses from the
three calibration models and exhibited mixed refusal behavior. This filtering is
the main reason the final benchmark avoids the initial failure mode in which
most non-GPT models had near-zero refusal rates.

\subsection{Fixed-Size Selection Ablation}
We also tested whether an additional fixed-size selection step was necessary
after calibration. A pilot comparison between a fixed 50-prompt subset and the
full calibrated pool showed that the model ranking and refusal-rate gradient
remained visible without compressing the dataset. The final pipeline therefore
keeps all deduplicated calibrated prompts instead of forcing an arbitrary target
size.

\section{Discussion}
\subsection{Why Calibration Is Needed}
Early pilot runs showed that safe code prompts could be too easy for many
models, leading to near-zero over-refusal outside a few GPT-style systems.
Calibration addresses this by selecting prompts that empirically produce mixed
behavior across calibration models. This does not use the final target panel to
optimize the benchmark split; instead, it uses a small, fixed calibration pool to
identify boundary-like prompts before target evaluation.

\subsection{Benchmark and Evaluation Decoupling}
The benchmark is fixed after 03C. Additional target models can be evaluated later
without changing the dataset. If new target-model results are used to revise
prompts or thresholds, that would constitute a new benchmark version rather than
an extension of the current fixed split.

\subsection{Limitations}
\benchmark currently focuses on six cybersecurity-adjacent code categories. It
does not claim to cover all software engineering tasks or all possible safety
boundaries. Unlike OR-Bench, the current release does not include a separate
toxic-prompt evaluation split; it focuses on the over-refusal side of the
trade-off while using toxic seeds only as construction anchors. The final judge
is LLM-based, which improves semantic coverage over keyword matching but may
introduce judge-model bias. Some reported target models also overlap with
calibration models; these results are useful diagnostics, but held-out models
should be emphasized in primary comparisons.

\subsection{Reproducibility}
The implementation is organized as a six-phase pipeline. Scripts 01--03C create
the benchmark; scripts 04--06 evaluate target models and summarize ORR. The
fixed dataset is stored as \texttt{dataset/full\_all/03c\_selected\_records.jsonl}.
The final metric report is stored as
\texttt{dataset/full\_all/06\_metrics\_report.md}. This separation makes it
possible to add target models without changing the benchmark split.

\section{Ethical Considerations}
The benchmark is designed to measure over-refusal on prompts that are safe to
fulfill. Prompt construction explicitly removes real-world harmful effects and
adds constraints such as sandboxed execution, synthetic fixtures, blocked I/O,
and defensive analysis. We avoid publishing operational harmful instructions in
the paper. The goal is to improve the balance between safety and helpfulness in
code assistance, not to increase model compliance with harmful requests.

\section{Conclusion}
We introduced \benchmark, a code-oriented over-refusal benchmark built from
controlled-risk rewrites and filtered through safety verification and
refusal-potential calibration. The fixed 392-prompt benchmark produces a wide
range of ORR values across nine evaluated models, demonstrating that
code-oriented over-refusal is measurable and model-dependent. Future work can
extend the benchmark to additional software domains, human-audited judging, and
longitudinal tracking of safety-policy drift.

\bibliographystyle{IEEEtran}
\bibliography{references}

\end{document}
